\section{Conclusion}\label{conclusion}

\subsection{Synthèse}\label{synthuxe8se}
Avant de réduire son impact environnental liés à l'IA générative, il faut comprendre les facteurs qui l'influencent.
On peut les regrouper en 3 catégories : le modèle, son hébergement, et les tokens.

Des outils existent pour mesurer celui des modèles ouverts avec précision et certitude.
Les modèles propriétaires les plus puissants comme ceux d'OpenAI, Google ou Anthropic sont des boîtes noires et ne permettent pas d'établir une mesure certaine. 
Des organismes indépendants développent des méthodes d'estimation basées sur de nombreuses hypothèses non vérifiables, mais sur lesquels on peut s'appuyer pour obtenir des ordres de grandeur.

La calculatrice que j'ai développé en parallèle de cet article permet à l'utilisateur final d'un chatbot de quantifier la réduction de son impact en mettant en place des bonnes pratiques : \url{https://felixmortas.com/ai-inf-calculator}

A son niveau, les trois principaux leviers qu'il peut actionner instantanément sont :
\begin{itemize}
    \item la priorisation d'une recherche web classique, en désactivant les recherches par IA automatiques.
    \item le choix d'un modèle de petite taille avec un effort faible, adapté à la tâche à réaliser
    \item la maîtrise de la quantité de tokens traités en entrée et générés en sortie.
\end{itemize}

En tant que développeur, l'amplitude d'action est plus large et touche des concepts invisibles pour l'utilisateur final. 

A tous les niveaux, il est possible de drastiquement réduire l'impact environnemental liée à son utilisation de l'IA, tout en augmentant la qualité des réponses produites.

\subsection{Ouverture}

Cet article est une synthèse des connaissances que j’ai accumulées au cours d’un an et demi de recherche, ainsi que de ma veille informationnelle continue sur le sujet.
Je tiens tout particulièrement à remercier Boris RUF, qui a partagé avec moi ses connaissances et ses compétences tout au long de cette période, alors que j'exerçais en tant que Data Scientist chez AXA Group Operations.
Je souhaite également remercier Céline LESCOP, qui m’a donné l’opportunité et l’envie d’approfondir le sujet du numérique écoresponsable en me permettant de participer au programme Digital Sustainability.

Si vous souhaitez approfondir le sujet, former vos équipes métier et IT aux bonnes pratiques, réaliser un bilan environnemental de vos outils existants ou concevoir de nouveaux outils, n’hésitez pas à me contacter. Je serai ravi d’intervenir au sein de votre organisation pour vous accompagner dans cette démarche.

Enfin, je reste ouvert à vos retours, critiques constructives et suggestions d’amélioration, tant sur le contenu de cet article que sur l’outil associé.
